\chapter{Groups}\label{groups}
\chapterstatus{complete}

\section{Introduction}\label{groups:section-introduction}

Groups describe symmetry. This chapter develops the basic theory: subgroups,
quotients and the isomorphism theorems, group actions and the Sylow theorems,
symmetric and alternating groups, products, solvable and nilpotent groups,
free groups and presentations, and the structure of finitely generated
abelian groups. Groups appear throughout the project: as fundamental groups
in topology (Chapter~\ref{homotopy}), as Galois groups
(Chapter~\ref{fields}), as structure groups of vector bundles, and as
Mumford--Tate groups of Hodge structures (Chapter~\ref{mumford-tate}).

References are \cite{lang-algebra}, \cite{rotman-groups} and
\cite{dummit-foote}; for free groups and amalgamated products see
\cite{serre-trees} and \cite{lyndon-schupp}.

\noindent\textbf{Prerequisites.} Chapter~\ref{sets} (Set Theory),
Chapter~\ref{number-systems} (Number Systems) and Chapter~\ref{categories}
(Categories).

\section{Groups and homomorphisms}\label{groups:section-groups}

\begin{definition}\label{groups:definition-group}
A \emph{group} is a set $G$ together with a map $G \times G \to G$, called the
\emph{group law} and written $(a, b) \mapsto a \cdot b$, and an element $e \in G$, such that
for all $a, b, c \in G$:
\begin{enumerate}
\item (associativity) $(a \cdot b) \cdot c = a \cdot (b \cdot c)$;
\item (neutral element) $e \cdot a = a \cdot e = a$;
\item (inverses) there is an element $a' \in G$ with $a \cdot a' = a' \cdot a = e$.
\end{enumerate}
The group is \emph{abelian} if $a \cdot b = b \cdot a$ for all $a, b \in G$. The \emph{order}
of $G$ is its cardinality $|G|$.
\end{definition}

\begin{notation}\label{groups:notation-multiplicative-additive}
Two notations are used for the group law.
\begin{itemize}
\item \emph{Multiplicative notation}, used for all groups: the group law is written
$a \cdot b$ or simply $ab$, the neutral element $e$ (or $1$), the inverse of $a$ is
written $a^{-1}$ (it is unique by Lemma~\ref{groups:lemma-basic}), and $a^n$ denotes
the $n$-th power (Lemma~\ref{groups:lemma-powers}).
\item \emph{Additive notation}, used only for abelian groups: the group law is
written $a + b$, the neutral element $0$, the inverse $-a$, and $na$ the $n$-th
multiple.
\end{itemize}
Every statement in multiplicative notation applies to abelian groups written
additively through the dictionary $ab \leftrightarrow a + b$, $e \leftrightarrow 0$,
$a^{-1} \leftrightarrow -a$, $a^n \leftrightarrow na$. General results in this chapter are
stated multiplicatively. For groups arising from a familiar structure we keep
the familiar symbol: $+$ for $(\ZZ, +)$, multiplication for $\QQ^\times$ or for
invertible matrices, and composition $\circ$ (written as juxtaposition) for groups of
bijections. The symbol $*$ never denotes a group law; $G_1 * G_2$ is the free
product of Definition~\ref{groups:definition-amalgam}. When the group law has to be
named, we write $(G, \cdot)$ or $(G, +)$.
\end{notation}

\begin{lemma}\label{groups:lemma-basic}
In a group the neutral element and inverses are unique, $(ab)^{-1} = b^{-1}a^{-1}$,
$(a^{-1})^{-1} = a$, and $ab = ac$ or $ba = ca$ implies $b = c$.
\end{lemma}

\begin{proof}
If $e'$ is also neutral, $e = ee' = e'$. If $b$ and $c$ are inverses of $a$,
then $b = b(ac) = (ba)c = c$; we write $a^{-1}$ for the inverse. The element
$b^{-1}a^{-1}$ is inverse to $ab$, and $a$ is inverse to $a^{-1}$. Cancellation
follows by multiplying with $a^{-1}$.
\end{proof}

\begin{example}\label{groups:example-groups}
\begin{enumerate}
\item $(\ZZ, +)$, $(\QQ, +)$, $(\RR, +)$ and $(\CC, +)$ are abelian groups, as are
$\QQ^\times = \QQ \setminus \{0\}$, $\RR^\times$ and $\CC^\times$ under multiplication.
More generally, for a commutative ring $R$ the invertible elements form a group
$R^\times$.
\item For a set $X$, the bijections $X \to X$ form the \emph{symmetric group}
$\mathrm{Sym}(X)$ under composition. We write $S_n = \mathrm{Sym}(\{1, \ldots, n\})$.
\item For a field $k$, the invertible $n \times n$ matrices form the group
$\GL_n(k)$ (Chapter~\ref{linear-algebra}).
\item The \emph{trivial group} has one element.
\item A group with one object as a category (Example~\ref{categories:example-monoid-as-category})
is a category with one object in which every morphism is an isomorphism.
\end{enumerate}
\end{example}

\begin{definition}\label{groups:definition-homomorphism}
A \emph{homomorphism} of groups is a map $f \colon G \to H$ with $f(ab) = f(a)f(b)$
for all $a, b \in G$. Its \emph{kernel} is $\ker f = \{a \in G : f(a) = e\}$ and its
\emph{image} is $\im f = f[G]$. An \emph{isomorphism} is a bijective homomorphism;
an \emph{automorphism} of $G$ is an isomorphism $G \to G$, and they form a group
$\Aut(G)$. Small groups and homomorphisms form a locally small category
$\cat{Grp}$, with the full subcategory $\cat{Ab}$ of abelian groups.
\end{definition}

\begin{lemma}\label{groups:lemma-homomorphism}
A homomorphism $f$ satisfies $f(e) = e$ and $f(a^{-1}) = f(a)^{-1}$; it is
injective if and only if $\ker f = \{e\}$; and the inverse of a bijective
homomorphism is a homomorphism.
\end{lemma}

\begin{proof}
$f(e) = f(e)f(e)$ gives $f(e) = e$ by cancellation, and then
$f(a)f(a^{-1}) = f(e) = e$. If $\ker f$ is trivial and $f(a) = f(b)$, then
$f(ab^{-1}) = e$, so $a = b$. If $f$ is bijective with inverse $g$, then
$f(g(x)g(y)) = xy$, so $g(xy) = g(x)g(y)$.
\end{proof}

\begin{definition}\label{groups:definition-subgroup}
A \emph{subgroup} of $G$ is a subset $H$ containing $e$ with $ab \in H$ and
$a^{-1} \in H$ for all $a, b \in H$; it is a group with the restricted operation.
We write $H \leq G$. For a subset $S \subseteq G$, the \emph{subgroup generated by}
$S$, written $\langle S \rangle$, is the intersection of all subgroups containing $S$.
\end{definition}

\begin{lemma}\label{groups:lemma-generated-subgroup}
The subgroup $\langle S \rangle$ consists of all finite products
$s_1^{\epsilon_1} \cdots s_k^{\epsilon_k}$ with $k \geq 0$, $s_i \in S$ and
$\epsilon_i \in \{1, -1\}$ (the empty product being $e$). Kernels and images of
homomorphisms are subgroups.
\end{lemma}

\begin{proof}
The set of such products contains $S$ and is closed under products and
inverses (Lemma~\ref{groups:lemma-basic}), so it contains $\langle S \rangle$; and
every subgroup containing $S$ contains all such products. The second statement
follows from Lemma~\ref{groups:lemma-homomorphism}.
\end{proof}

\begin{lemma}\label{groups:lemma-powers}
Let $G$ be a group and $g \in G$. There is a unique homomorphism
$\ZZ \to G$ sending $1$ to $g$; we write $n \mapsto g^n$. Thus
$g^{m+n} = g^m g^n$ and $(g^m)^n = g^{mn}$ for all $m, n \in \ZZ$.
\end{lemma}

\begin{proof}
Uniqueness is clear since $1$ generates $\ZZ$. For existence define $g^n$ for
$n \in \NN$ by recursion ($g^0 = e$, $g^{n+1} = g^n g$) and put
$g^{-n} = (g^{-1})^n$. One checks $g^{m+1} = g^m g$ for all $m \in \ZZ$ (for
$m = -n - 1 < 0$ this reads $(g^{-1})^{n} = (g^{-1})^{n+1} g$, which holds since
$g^{-1}g = e$), and then $g^{m+n} = g^m g^n$ follows by induction on $n \geq 0$ and
on $-n \geq 0$. The last formula holds because $n \mapsto g^{mn}$ and
$n \mapsto (g^m)^n$ are both homomorphisms sending $1$ to $g^m$.
\end{proof}

\begin{proposition}\label{groups:proposition-subgroups-integers}
Every subgroup of $\ZZ$ is of the form $n\ZZ = \{nk : k \in \ZZ\}$ for a unique
$n \in \NN$.
\end{proposition}

\begin{proof}
Let $H \leq \ZZ$. If $H = \{0\}$ take $n = 0$. Otherwise $H$ contains a positive
element, and we let $n$ be the least one. Then $n\ZZ \subseteq H$, and for
$h \in H$, division with remainder (Theorem~\ref{number-systems:theorem-division})
gives $h = qn + r$ with $0 \leq r < n$ and $r = h - qn \in H$, so $r = 0$. The number
$n$ is the least positive element of $n\ZZ$, hence unique.
\end{proof}

\begin{definition}\label{groups:definition-order-element}
The \emph{order} of $g \in G$ is the least $n \geq 1$ with $g^n = e$ if it exists, and
$\infty$ otherwise. By Proposition~\ref{groups:proposition-subgroups-integers},
the kernel of $n \mapsto g^n$ is $\operatorname{ord}(g)\ZZ$ (with $\infty \ZZ := \{0\}$).
A group is \emph{cyclic} if it is generated by one element.
\end{definition}

\section{Cosets and Lagrange's theorem}\label{groups:section-cosets}

\begin{definition}\label{groups:definition-cosets}
Let $H \leq G$. A \emph{left coset} of $H$ is a set $gH = \{gh : h \in H\}$, and a
right coset is a set $Hg$. The set of left cosets is $G/H$, and the \emph{index}
of $H$ is $[G : H] = |G/H|$.
\end{definition}

\begin{lemma}\label{groups:lemma-cosets-partition}
The left cosets of $H$ are the classes of the equivalence relation
$a \sim b \iff a^{-1}b \in H$; hence they partition $G$. Each left coset is in
bijection with $H$, and $gH \mapsto Hg^{-1}$ is a bijection from left to right
cosets.
\end{lemma}

\begin{proof}
The relation is an equivalence relation because $H$ is a subgroup, and the class
of $a$ is $aH$. The map $h \mapsto gh$ is a bijection $H \to gH$. Finally $aH = bH$
iff $a^{-1}b \in H$ iff $(a^{-1}b)^{-1} = b^{-1}a \in H$ iff $Ha^{-1} = Hb^{-1}$.
\end{proof}

\begin{theorem}[Lagrange]\label{groups:theorem-lagrange}
Let $G$ be finite and $H \leq G$. Then $|G| = [G : H] \cdot |H|$. In particular the
order of $H$ divides the order of $G$, and the order of every element divides
$|G|$, so $g^{|G|} = e$ for all $g \in G$.
\end{theorem}

\begin{proof}
By Lemma~\ref{groups:lemma-cosets-partition}, $G$ is the disjoint union of
$[G:H]$ sets of size $|H|$. The subgroup $\langle g \rangle$ has order
$\operatorname{ord}(g)$, since $g^0, \ldots, g^{\operatorname{ord}(g)-1}$ are
distinct and every power $g^n$ equals $g^r$ with $r$ the remainder of $n$.
\end{proof}

\begin{corollary}\label{groups:corollary-prime-order}
A group of prime order $p$ is cyclic, generated by any element $\neq e$.
\end{corollary}

\begin{proof}
If $g \neq e$, then $\operatorname{ord}(g) > 1$ divides $p$, so it equals $p$, and
$\langle g \rangle = G$.
\end{proof}

\begin{proposition}\label{groups:proposition-index-multiplicative}
If $K \leq H \leq G$ and the indices are finite, then $[G : K] = [G : H][H : K]$.
\end{proposition}

\begin{proof}
Let $(g_i)$ represent the left cosets of $H$ in $G$ and $(h_j)$ those of $K$ in $H$.
Then the $g_i h_j K$ are the left cosets of $K$ in $G$, and they are distinct:
$G = \bigsqcup_i g_i H = \bigsqcup_i \bigsqcup_j g_i h_j K$.
\end{proof}

\section{Normal subgroups and quotients}\label{groups:section-quotients}

\begin{definition}\label{groups:definition-normal}
A subgroup $N \leq G$ is \emph{normal}, written $N \trianglelefteq G$, if
$gNg^{-1} = N$ for all $g \in G$, equivalently $gN = Ng$ for all $g$. A group
$G \neq \{e\}$ is \emph{simple} if its only normal subgroups are $\{e\}$ and $G$.
\end{definition}

\begin{lemma}\label{groups:lemma-normal-examples}
Kernels of homomorphisms are normal. Every subgroup of an abelian group is
normal. A subgroup of index $2$ is normal.
\end{lemma}

\begin{proof}
If $f(n) = e$ then $f(gng^{-1}) = f(g)f(g)^{-1} = e$. The second claim is clear.
If $[G:H] = 2$, then for $g \notin H$ both $gH$ and $Hg$ are the complement of $H$.
\end{proof}

\begin{theorem}\label{groups:theorem-quotient}
Let $N \trianglelefteq G$. The rule $(aN)(bN) = abN$ is a well-defined operation
making $G/N$ a group, the \emph{quotient group}, and the projection
$\pi \colon G \to G/N$ is a surjective homomorphism with kernel $N$. If
$f \colon G \to H$ is a homomorphism with $N \subseteq \ker f$, there is a unique
homomorphism $\bar f \colon G/N \to H$ with $\bar f \circ \pi = f$.
\end{theorem}

\begin{proof}
If $aN = a'N$ and $bN = b'N$, write $a' = an$, $b' = bm$. Then
$a'b' = anbm = ab(b^{-1}nb)m \in abN$ since $N$ is normal. The group axioms follow
from those of $G$, with neutral element $N$ and $(aN)^{-1} = a^{-1}N$. Clearly $\pi$
is a surjective homomorphism, and $\pi(a) = N$ iff $a \in N$. If $N \subseteq \ker f$
then $f$ is constant on cosets, and $\bar f$ exists and is unique by
Proposition~\ref{sets:proposition-quotient-property}; it is a homomorphism because
$\pi$ is surjective.
\end{proof}

\begin{theorem}[Isomorphism theorems]\label{groups:theorem-isomorphism-theorems}
\begin{enumerate}
\item If $f \colon G \to H$ is a homomorphism, $\bar f \colon G/\ker f \to \im f$ is an
isomorphism.
\item If $H \leq G$ and $N \trianglelefteq G$, then $HN$ is a subgroup, $N \trianglelefteq HN$,
$H \cap N \trianglelefteq H$, and $H/(H \cap N) \cong HN/N$.
\item If $N \leq M$ are normal subgroups of $G$, then $M/N \trianglelefteq G/N$ and
$(G/N)/(M/N) \cong G/M$.
\item (Correspondence theorem) If $N \trianglelefteq G$, then $K \mapsto K/N$ is a bijection
from the set of subgroups of $G$ containing $N$ to the set of subgroups of $G/N$,
preserving inclusions, normality and indices.
\end{enumerate}
\end{theorem}

\begin{proof}
(1) $\bar f$ is surjective onto $\im f$ and has trivial kernel.
(2) $HN = NH$ because $hN = Nh$, so $HN$ is closed under products and inverses.
The composite $H \to HN \to HN/N$ is surjective, since $hnN = hN$, with kernel
$H \cap N$; apply (1).
(3) The map $G/N \to G/M$, $gN \mapsto gM$, is a well-defined surjective homomorphism
(Theorem~\ref{groups:theorem-quotient}) with kernel $M/N$; apply (1).
(4) The inverse map sends $\bar K \leq G/N$ to $\pi^{-1}[\bar K]$; since $\pi$ is
surjective, $\pi[\pi^{-1}[\bar K]] = \bar K$, and $\pi^{-1}[\pi[K]] = KN = K$ when
$N \subseteq K$. Normality is preserved because $\pi$ is surjective, and
$gK \mapsto (gN)(K/N)$ is a bijection $G/K \to (G/N)/(K/N)$.
\end{proof}

\begin{example}\label{groups:example-cyclic-groups}
For $n \geq 1$, the quotient $\ZZ/n\ZZ$ is a cyclic group of order $n$, with
elements the classes of $0, 1, \ldots, n-1$. If $G = \langle g \rangle$ is cyclic, the
homomorphism $\ZZ \to G$, $k \mapsto g^k$ is surjective with kernel
$\operatorname{ord}(g)\ZZ$, so $G \cong \ZZ$ if $g$ has infinite order and
$G \cong \ZZ/n\ZZ$ if $\operatorname{ord}(g) = n$. Thus a cyclic group is determined
up to isomorphism by its order. By the correspondence theorem the subgroups of
$\ZZ/n\ZZ$ are the $d\ZZ/n\ZZ$ for $d \mid n$, one for each divisor.
\end{example}

\begin{definition}\label{groups:definition-center-conjugation}
The \emph{center} of $G$ is $Z(G) = \{z \in G : zg = gz \text{ for all } g\}$. For
$g \in G$, \emph{conjugation} $c_g(x) = gxg^{-1}$ is an automorphism; automorphisms
of this form are \emph{inner}. The \emph{centraliser} of $x$ is
$C_G(x) = \{g : gx = xg\}$ and the \emph{normaliser} of $H \leq G$ is
$N_G(H) = \{g : gHg^{-1} = H\}$.
\end{definition}

\begin{lemma}\label{groups:lemma-center-cyclic}
$Z(G)$ is a normal subgroup, and $g \mapsto c_g$ is a homomorphism
$G \to \Aut(G)$ with kernel $Z(G)$. If $G/Z(G)$ is cyclic, then $G$ is abelian.
\end{lemma}

\begin{proof}
$c_{gh} = c_g \circ c_h$, and $c_g = \id$ iff $g \in Z(G)$. If $G/Z(G)$ is generated
by $gZ(G)$, every element of $G$ is $g^k z$ with $z \in Z(G)$, and such elements
commute with each other.
\end{proof}

\section{Products}\label{groups:section-products}

\begin{proposition}\label{groups:proposition-products}
For a family of groups $(G_i)_{i \in I}$, the set $\prod_i G_i$ with the componentwise
operation is a group, and with the projections it is a product in $\cat{Grp}$
(Example~\ref{categories:example-limits}). Let $N, M \trianglelefteq G$ with
$N \cap M = \{e\}$ and $NM = G$. Then $nm = mn$ for $n \in N$, $m \in M$, and
$N \times M \to G$, $(n, m) \mapsto nm$, is an isomorphism.
\end{proposition}

\begin{proof}
The first statement is clear. For $n \in N$ and $m \in M$, the commutator
$nmn^{-1}m^{-1}$ lies in $N$ (as $mn^{-1}m^{-1} \in N$) and in $M$ (as
$nmn^{-1} \in M$), hence is $e$. So the map is a homomorphism; it is surjective
since $NM = G$, and injective since $nm = e$ implies $n = m^{-1} \in N \cap M$.
\end{proof}

\begin{definition}\label{groups:definition-semidirect}
Let $N$ and $H$ be groups and $\varphi \colon H \to \Aut(N)$, $h \mapsto \varphi_h$, a
homomorphism. The \emph{semidirect product} $N \rtimes_\varphi H$ is the set
$N \times H$ with the operation $(n, h)(n', h') = (n\varphi_h(n'), hh')$.
\end{definition}

\begin{proposition}\label{groups:proposition-semidirect}
$N \rtimes_\varphi H$ is a group, in which $N \times \{e\}$ is a normal subgroup
isomorphic to $N$, $\{e\} \times H$ is a subgroup isomorphic to $H$, and
$(e, h)(n, e)(e, h)^{-1} = (\varphi_h(n), e)$. Conversely, if $G$ is a group with
$N \trianglelefteq G$ and $H \leq G$ such that $N \cap H = \{e\}$ and $NH = G$, then
$(n, h) \mapsto nh$ is an isomorphism $N \rtimes_\varphi H \to G$, where
$\varphi_h(n) = hnh^{-1}$.
\end{proposition}

\begin{proof}
Associativity:
$((n,h)(n',h'))(n'',h'') = (n\varphi_h(n')\varphi_{hh'}(n''), hh'h'')$ and
$(n,h)((n',h')(n'',h'')) = (n\varphi_h(n'\varphi_{h'}(n'')), hh'h'')$ agree since
$\varphi$ is a homomorphism into $\Aut(N)$. The neutral element is $(e, e)$ and
$(n,h)^{-1} = (\varphi_{h^{-1}}(n^{-1}), h^{-1})$. The projection
$(n, h) \mapsto h$ is a homomorphism with kernel $N \times \{e\}$, and the conjugation
formula is a direct computation. For the converse, every $g \in G$ is uniquely
$nh$ (if $nh = n'h'$ then $n'^{-1}n = h'h^{-1} \in N \cap H$), and
$nh \cdot n'h' = n(hn'h^{-1})hh'$, which is the product rule.
\end{proof}

\begin{example}\label{groups:example-dihedral}
For $n \geq 3$ let $\varphi \colon \ZZ/2\ZZ \to \Aut(\ZZ/n\ZZ)$ send the generator to
$x \mapsto -x$. The group $D_n = \ZZ/n\ZZ \rtimes_\varphi \ZZ/2\ZZ$ of order $2n$ is the
\emph{dihedral group}; it is the group of symmetries of a regular $n$-gon (rotations
and reflections). It is not abelian, since $(0,1)(1,0) = (-1, 1) \neq (1, 1) = (1,0)(0,1)$.
\end{example}

\section{Group actions}\label{groups:section-actions}

\begin{definition}\label{groups:definition-action}
An \emph{action} of a group $G$ on a set $X$ is a map $G \times X \to X$,
$(g, x) \mapsto gx$, with $ex = x$ and $g(hx) = (gh)x$; equivalently a homomorphism
$G \to \mathrm{Sym}(X)$. The \emph{orbit} of $x$ is $Gx = \{gx : g \in G\}$, its
\emph{stabiliser} is $G_x = \{g : gx = x\}$, and $X^G = \{x : gx = x \text{ for all } g\}$
is the set of \emph{fixed points}. The action is \emph{transitive} if there is
only one orbit, and \emph{free} if all stabilisers are trivial.
\end{definition}

\begin{theorem}[Orbit--stabiliser]\label{groups:theorem-orbit-stabiliser}
For an action of $G$ on $X$ and $x \in X$, the stabiliser $G_x$ is a subgroup and
$gG_x \mapsto gx$ is a bijection $G/G_x \to Gx$. In particular
$|Gx| = [G : G_x]$, and the orbits partition $X$.
\end{theorem}

\begin{proof}
$gG_x = hG_x$ iff $h^{-1}g \in G_x$ iff $gx = hx$, so the map is well defined and
injective; it is surjective by definition. The relation ``$y \in Gx$'' is an
equivalence relation since $G$ is a group.
\end{proof}

\begin{example}\label{groups:example-actions}
\begin{enumerate}
\item $G$ acts on itself by left multiplication, freely and transitively. This
gives an injective homomorphism $G \to \mathrm{Sym}(G)$ (\emph{Cayley's theorem}):
every group is isomorphic to a group of permutations.
\item $G$ acts on itself by conjugation; the orbits are the \emph{conjugacy
classes} and $G_x = C_G(x)$.
\item $G$ acts on the set of its subgroups by conjugation, with stabiliser
$N_G(H)$; so the number of conjugates of $H$ is $[G : N_G(H)]$.
\item If $H \leq G$, then $G$ acts transitively on $G/H$ by $g(aH) = gaH$.
\end{enumerate}
\end{example}

\begin{corollary}[Class equation]\label{groups:corollary-class-equation}
Let $G$ be finite and let $x_1, \ldots, x_r$ represent the conjugacy classes with
more than one element. Then $|G| = |Z(G)| + \sum_{i=1}^r [G : C_G(x_i)]$.
\end{corollary}

\begin{proof}
The conjugacy classes partition $G$; the one-element classes are the elements of
$Z(G)$, and the class of $x_i$ has $[G : C_G(x_i)]$ elements by
Theorem~\ref{groups:theorem-orbit-stabiliser}.
\end{proof}

\begin{definition}\label{groups:definition-p-group}
Let $p$ be a prime. A \emph{$p$-group} is a finite group whose order is a power
of $p$.
\end{definition}

\begin{lemma}\label{groups:lemma-p-group-fixed-points}
Let a $p$-group $P$ act on a finite set $X$. Then $|X| \equiv |X^P| \pmod p$.
\end{lemma}

\begin{proof}
Each orbit that is not a fixed point has size $[P : P_x] > 1$, a divisor of $|P|$,
hence divisible by $p$.
\end{proof}

\begin{proposition}\label{groups:proposition-p-group-center}
A nontrivial $p$-group has nontrivial center. Every group of order $p^2$ is
abelian.
\end{proposition}

\begin{proof}
Apply Lemma~\ref{groups:lemma-p-group-fixed-points} to the conjugation action:
$|Z(G)| \equiv |G| \equiv 0 \pmod p$ and $e \in Z(G)$, so $|Z(G)| \geq p$. If
$|G| = p^2$, then $|G/Z(G)| \in \{1, p\}$, so $G/Z(G)$ is cyclic
(Corollary~\ref{groups:corollary-prime-order}), and $G$ is abelian by
Lemma~\ref{groups:lemma-center-cyclic}.
\end{proof}

\begin{theorem}[Cauchy]\label{groups:theorem-cauchy}
If a prime $p$ divides the order of a finite group $G$, then $G$ has an element of
order $p$.
\end{theorem}

\begin{proof}
Let $X = \{(g_1, \ldots, g_p) \in G^p : g_1 g_2 \cdots g_p = e\}$. Since $g_p$ is
determined by the others, $|X| = |G|^{p-1} \equiv 0 \pmod p$. The group $\ZZ/p\ZZ$
acts on $X$ by cyclic rotation: if $g_1 \cdots g_p = e$ then
$g_2 \cdots g_p g_1 = g_1^{-1}(g_1 \cdots g_p)g_1 = e$. The fixed points are the tuples
$(g, \ldots, g)$ with $g^p = e$. By Lemma~\ref{groups:lemma-p-group-fixed-points}
their number is divisible by $p$, and $(e, \ldots, e)$ is one of them, so there is
a fixed point with $g \neq e$.
\end{proof}

\section{The Sylow theorems}\label{groups:section-sylow}

\begin{definition}\label{groups:definition-sylow}
Let $G$ be finite and $|G| = p^a m$ with $p$ prime and $p \nmid m$. A \emph{Sylow
$p$-subgroup} of $G$ is a subgroup of order $p^a$. We write $\mathrm{Syl}_p(G)$ for
the set of them and $n_p = |\mathrm{Syl}_p(G)|$.
\end{definition}

\begin{theorem}[Sylow]\label{groups:theorem-sylow}
Let $G$ be a finite group, $p$ a prime and $|G| = p^a m$ with $p \nmid m$.
\begin{enumerate}
\item Sylow $p$-subgroups exist.
\item Every $p$-subgroup of $G$ is contained in a Sylow $p$-subgroup, and any two
Sylow $p$-subgroups are conjugate.
\item $n_p$ divides $m$ and $n_p \equiv 1 \pmod p$.
\end{enumerate}
\end{theorem}

\begin{proof}
(1) By induction on $|G|$; if $a = 0$ there is nothing to prove. Suppose
$p \mid |Z(G)|$. By Theorem~\ref{groups:theorem-cauchy} there is $z \in Z(G)$ of order
$p$; $N = \langle z \rangle$ is normal, and $G/N$ has order $p^{a-1}m$, so by
induction it has a subgroup of order $p^{a-1}$, whose preimage in $G$ has order
$p^a$. Suppose $p \nmid |Z(G)|$. In the class equation
(Corollary~\ref{groups:corollary-class-equation}) the left side is divisible by
$p$, so some $[G : C_G(x_i)]$ is not divisible by $p$. Then $C_G(x_i)$ is a proper
subgroup of order divisible by $p^a$, and by induction it contains a subgroup of
order $p^a$.

(2) Let $P$ be a Sylow $p$-subgroup and $Q$ a $p$-subgroup. Let $Q$ act on $G/P$
by left multiplication. Since $|G/P| = m \not\equiv 0 \pmod p$,
Lemma~\ref{groups:lemma-p-group-fixed-points} gives a fixed point $gP$. Then
$QgP = gP$, so $g^{-1}Qg \subseteq P$ and $Q \subseteq gPg^{-1}$, which is a Sylow
$p$-subgroup. If $Q$ is itself Sylow, then $Q = gPg^{-1}$ by comparing orders.

(3) By (2), $G$ acts transitively by conjugation on $\mathrm{Syl}_p(G)$, so
$n_p = [G : N_G(P)]$, which divides $[G : P] = m$ since $P \leq N_G(P)$
(Proposition~\ref{groups:proposition-index-multiplicative}). Let $P$ act on
$\mathrm{Syl}_p(G)$ by conjugation. If $Q$ is a fixed point, then $P \leq N_G(Q)$, so
$P$ and $Q$ are Sylow $p$-subgroups of $N_G(Q)$, hence conjugate in $N_G(Q)$ by
(2); as $Q$ is normal in $N_G(Q)$, $P = Q$. So $P$ is the only fixed point, and
$n_p \equiv 1 \pmod p$ by Lemma~\ref{groups:lemma-p-group-fixed-points}.
\end{proof}

\begin{corollary}\label{groups:corollary-unique-sylow}
A Sylow $p$-subgroup is normal if and only if $n_p = 1$.
\end{corollary}

\begin{proof}
By Theorem~\ref{groups:theorem-sylow}(2) the Sylow $p$-subgroups are the
conjugates of one of them.
\end{proof}

\begin{example}\label{groups:example-order-15}
Every group $G$ of order $15$ is cyclic. Indeed $n_3$ divides $5$ and is $1$
modulo $3$, so $n_3 = 1$; similarly $n_5 = 1$. The Sylow subgroups $P_3$, $P_5$ are
normal, cyclic of prime order, and $P_3 \cap P_5 = \{e\}$ by Lagrange. Then
$|P_3P_5| = 15$ (the map $P_3 \times P_5 \to G$ is injective), so
$G \cong \ZZ/3\ZZ \times \ZZ/5\ZZ$ by Proposition~\ref{groups:proposition-products},
and the element $(1, 1)$ has order $15$.
\end{example}

\section{Symmetric and alternating groups}\label{groups:section-symmetric}

\begin{definition}\label{groups:definition-cycles}
Let $n \geq 1$. For distinct $a_1, \ldots, a_k \in \{1, \ldots, n\}$, the \emph{cycle}
$(a_1\, a_2\, \cdots\, a_k) \in S_n$ sends $a_i$ to $a_{i+1}$ for $i < k$, $a_k$ to $a_1$,
and fixes all other points. A cycle of length $2$ is a \emph{transposition}. Two
cycles are \emph{disjoint} if they move disjoint sets of points. Products in $S_n$
are compositions, applied from right to left.
\end{definition}

\begin{proposition}\label{groups:proposition-cycle-decomposition}
Every $\sigma \in S_n$ is a product of disjoint cycles of length $\geq 2$, uniquely up
to order; disjoint cycles commute. Every element of $S_n$ is a product of
transpositions. For $\tau \in S_n$,
$\tau (a_1\, \cdots\, a_k) \tau^{-1} = (\tau(a_1)\, \cdots\, \tau(a_k))$, so two
permutations are conjugate if and only if they have the same cycle lengths.
\end{proposition}

\begin{proof}
The orbits of $\langle \sigma \rangle$ on $\{1, \ldots, n\}$ partition it; on an orbit
$\{a, \sigma(a), \ldots, \sigma^{k-1}(a)\}$ of size $k$, $\sigma$ acts as the cycle
$(a\, \sigma(a)\, \cdots\, \sigma^{k-1}(a))$, and $\sigma$ is the product of these cycles for
the orbits of size $\geq 2$. Conversely the cycles in any such decomposition must be
the orbits, which gives uniqueness. Disjoint cycles commute since they act on
disjoint sets. We have $(a_1\, \cdots\, a_k) = (a_1\, a_k)(a_1\, a_{k-1}) \cdots (a_1\, a_2)$,
as one checks by evaluating both sides. The conjugation formula follows by
evaluating both sides at $\tau(a_i)$ and at points not of this form. If $\sigma$ and
$\sigma'$ have the same cycle lengths, a bijection matching their cycles gives $\tau$
with $\tau\sigma\tau^{-1} = \sigma'$.
\end{proof}

\begin{proposition}\label{groups:proposition-sign}
There is a unique homomorphism $\operatorname{sgn} \colon S_n \to \{1, -1\}$ sending
every transposition to $-1$ (for $n \geq 2$). It is given by
$\operatorname{sgn}(\sigma) = \prod_{1 \leq i < j \leq n} \frac{\sigma(j) - \sigma(i)}{j - i}$.
\end{proposition}

\begin{proof}
For $\sigma \in S_n$ let $s(\sigma)$ be the product in the statement, a rational
number. The factors $\sigma(j) - \sigma(i)$ run, up to sign, over all differences
$j - i$ with $i < j$, so $s(\sigma) = \pm 1$. For $\sigma, \tau \in S_n$,
\[
s(\sigma\tau) = \prod_{i < j} \frac{\sigma(\tau(j)) - \sigma(\tau(i))}{\tau(j) - \tau(i)} \cdot
\prod_{i < j} \frac{\tau(j) - \tau(i)}{j - i}.
\]
The expression $\frac{\sigma(b) - \sigma(a)}{b - a}$ is symmetric in $a$ and $b$, and
$\{\tau(i), \tau(j)\}$ runs over all two-element subsets as $\{i, j\}$ does; so the first
product is $s(\sigma)$, and $s$ is a homomorphism. For the transposition $(1\,2)$, the
only factor that is negative is the one with $(i, j) = (1, 2)$, so $s((1\,2)) = -1$.
All transpositions are conjugate by Proposition~\ref{groups:proposition-cycle-decomposition},
hence have sign $-1$. Uniqueness holds because transpositions generate $S_n$.
\end{proof}

\begin{definition}\label{groups:definition-alternating}
The \emph{alternating group} $A_n$ is the kernel of $\operatorname{sgn}$. Permutations in
$A_n$ are \emph{even}. For $n \geq 2$, $[S_n : A_n] = 2$.
\end{definition}

\begin{lemma}\label{groups:lemma-three-cycles}
For $n \geq 3$, the group $A_n$ is generated by the $3$-cycles. For $n \geq 5$, any two
$3$-cycles are conjugate in $A_n$.
\end{lemma}

\begin{proof}
A $3$-cycle is even, being a product of two transpositions. Every even permutation
is a product of an even number of transpositions, and for distinct $a, b, c, d$ we
have $(a\,b)(a\,c) = (a\,c\,b)$ and $(a\,b)(c\,d) = (a\,b\,c)(b\,c\,d)$, as one checks by
evaluation. For the second claim, two $3$-cycles $(a\,b\,c)$ and $(a'\,b'\,c')$ are
conjugate by some $\tau \in S_n$. If $\tau$ is odd, choose $d, e \notin \{a, b, c\}$
distinct (possible since $n \geq 5$); then $\tau(d\,e)$ is even and also conjugates
$(a\,b\,c)$ to $(a'\,b'\,c')$, because $(d\,e)$ commutes with $(a\,b\,c)$.
\end{proof}

\begin{theorem}\label{groups:theorem-alternating-simple}
For $n \geq 5$ the group $A_n$ is simple.
\end{theorem}

\begin{proof}
Let $N \trianglelefteq A_n$ with $N \neq \{e\}$. By Lemma~\ref{groups:lemma-three-cycles}
it suffices to show that $N$ contains a $3$-cycle. Choose $\sigma \in N \setminus \{e\}$
with the largest possible number of fixed points. For $\rho \in A_n$, the element
$\tau = (\rho\sigma\rho^{-1})\sigma^{-1}$ lies in $N$.

Suppose $\sigma$ moves exactly $4$ points. Being even, it is then of the form
$(a\,b)(c\,d)$. Choose $e \notin \{a, b, c, d\}$ and $\rho = (c\,d\,e)$. Then
$\rho\sigma\rho^{-1} = (a\,b)(d\,e)$ and $\tau = (a\,b)(d\,e)(a\,b)(c\,d) = (d\,e)(c\,d) = (c\,e\,d)$,
a $3$-cycle in $N$.

Suppose $\sigma$ moves at least $5$ points. If $\sigma$ has a cycle of length
$\geq 3$, write it as $(a\,b\,c\,\cdots)$ and choose further points $d, e$ moved by $\sigma$,
distinct from $a, b, c$. Otherwise $\sigma$ is a product of at least three disjoint
transpositions, and we write $\sigma = (a\,b)(c\,d)(e\,f)\cdots$. In both cases let
$\rho = (c\,d\,e)$. Every point fixed by $\sigma$ is fixed by $\rho$, hence by $\tau$; and
$\tau(\sigma(x)) = \rho\sigma\rho^{-1}(x)$. In the first case $\tau(b) = \tau(\sigma(a)) = \rho(\sigma(a)) = \rho(b) = b$,
and $\tau(c) = \tau(\sigma(b)) = \rho(\sigma(b)) = \rho(c) = d \neq c$. In the second case
$\tau(b) = \rho\sigma(a) = b$, $\tau(a) = \rho\sigma(b) = a$, and $\tau(d) = \tau(\sigma(c)) = \rho\sigma\rho^{-1}(c) = \rho(\sigma(e)) = \rho(f) = f \neq d$.
In both cases $\tau \neq e$, $\tau$ fixes all points fixed by $\sigma$ and at least one
point moved by $\sigma$, contradicting the choice of $\sigma$. So this case does not
occur. The remaining possibility is that $\sigma$ moves exactly $3$ points (an even
permutation cannot move exactly $2$), and then $\sigma$ is a $3$-cycle. In all cases $N$
contains a $3$-cycle.
\end{proof}

\begin{counterexample}\label{groups:counterexample-converse-lagrange}
The converse of Lagrange's theorem is false: $A_4$ has order $12$ but no subgroup
of order $6$. Such a subgroup $H$ would have index $2$, so $g^2 \in H$ for all
$g \in A_4$ (as $G/H$ has order $2$). But every $3$-cycle is a square:
$(a\,b\,c) = (a\,c\,b)^2$. So $H$ would contain the eight $3$-cycles of $A_4$.
\end{counterexample}

\begin{counterexample}\label{groups:counterexample-normal-not-transitive}
Normality is not transitive. In $A_4$, the subgroup
$V = \{e, (1\,2)(3\,4), (1\,3)(2\,4), (1\,4)(2\,3)\}$ is normal, being a union of
conjugacy classes by Proposition~\ref{groups:proposition-cycle-decomposition}. Its
subgroup $K = \{e, (1\,2)(3\,4)\}$ is normal in the abelian group $V$. But
$(1\,2\,3)(1\,2)(3\,4)(1\,2\,3)^{-1} = (2\,3)(1\,4) \notin K$, so $K$ is not normal in $A_4$.
\end{counterexample}

\section{Solvable and nilpotent groups}\label{groups:section-solvable}

\begin{definition}\label{groups:definition-commutator}
The \emph{commutator} of $a, b \in G$ is $[a, b] = aba^{-1}b^{-1}$. For subgroups
$H, K$, $[H, K]$ is the subgroup generated by the $[h, k]$. The \emph{derived
subgroup} is $G' = [G, G]$, and the \emph{derived series} is $G^{(0)} = G$,
$G^{(i+1)} = [G^{(i)}, G^{(i)}]$. The group $G$ is \emph{solvable} if
$G^{(n)} = \{e\}$ for some $n$.
\end{definition}

\begin{proposition}\label{groups:proposition-abelianization}
$[G, G]$ is a normal subgroup, $G^{\mathrm{ab}} = G/[G, G]$ is abelian, and a normal
subgroup $N$ has abelian quotient $G/N$ if and only if $[G, G] \subseteq N$. Every
homomorphism from $G$ to an abelian group factors uniquely through $G^{\mathrm{ab}}$;
thus $G \mapsto G^{\mathrm{ab}}$ is left adjoint to the inclusion $\cat{Ab} \to \cat{Grp}$.
\end{proposition}

\begin{proof}
$g[a, b]g^{-1} = [gag^{-1}, gbg^{-1}]$, so $[G, G]$ is normal. The quotient $G/N$ is
abelian iff $abN = baN$ for all $a, b$, iff all commutators lie in $N$. A
homomorphism $f$ to an abelian group kills all commutators, so it factors
uniquely by Theorem~\ref{groups:theorem-quotient}. The last statement is the
resulting natural bijection $\Hom_{\cat{Grp}}(G, A) \cong \Hom_{\cat{Ab}}(G^{\mathrm{ab}}, A)$.
\end{proof}

\begin{proposition}\label{groups:proposition-solvable}
\begin{enumerate}
\item $G$ is solvable if and only if there is a chain
$\{e\} = G_n \trianglelefteq G_{n-1} \trianglelefteq \cdots \trianglelefteq G_0 = G$ with all $G_{i}/G_{i+1}$ abelian.
\item Subgroups and quotients of solvable groups are solvable. If
$N \trianglelefteq G$ and both $N$ and $G/N$ are solvable, then $G$ is solvable.
\item $p$-groups are solvable, and $S_n$ is not solvable for $n \geq 5$.
\end{enumerate}
\end{proposition}

\begin{proof}
(1) The derived series is such a chain. Conversely, given a chain, induction shows
$G^{(i)} \subseteq G_i$, because $G_i/G_{i+1}$ abelian gives
$[G_i, G_i] \subseteq G_{i+1}$ (Proposition~\ref{groups:proposition-abelianization}).
(2) For $H \leq G$, $H^{(i)} \subseteq G^{(i)}$. For a surjection $f \colon G \to Q$,
$f[G^{(i)}] = Q^{(i)}$. If $(G/N)^{(n)}$ is trivial then $G^{(n)} \subseteq N$, and then
$G^{(n + m)} \subseteq N^{(m)}$ is trivial for large $m$.
(3) By induction on $|G|$, using (2) and
Proposition~\ref{groups:proposition-p-group-center}: $Z(G)$ is abelian and nontrivial,
and $G/Z(G)$ is a smaller $p$-group. For $n \geq 5$, $A_n$ is simple and not abelian,
so $[A_n, A_n] = A_n$ and $A_n$, hence $S_n$, is not solvable.
\end{proof}

\begin{example}\label{groups:example-s4-solvable}
$S_4$ is solvable: $\{e\} \trianglelefteq V \trianglelefteq A_4 \trianglelefteq S_4$, with $V$ as in
Counterexample~\ref{groups:counterexample-normal-not-transitive}, has quotients of
orders $4$, $3$ and $2$, which are abelian by
Proposition~\ref{groups:proposition-p-group-center}. The failure of solvability of
$S_n$ for $n \geq 5$ is the group-theoretic reason why polynomial equations of
degree $\geq 5$ cannot in general be solved by radicals (Chapter~\ref{fields}).
\end{example}

\begin{definition}\label{groups:definition-nilpotent}
The \emph{upper central series} of $G$ is defined by $Z_0 = \{e\}$ and
$Z_{i+1}/Z_i = Z(G/Z_i)$ (using the correspondence theorem). The group is
\emph{nilpotent} if $Z_n = G$ for some $n$.
\end{definition}

\begin{proposition}\label{groups:proposition-nilpotent}
$p$-groups are nilpotent, and nilpotent groups are solvable. The group $S_3$ is
solvable but not nilpotent.
\end{proposition}

\begin{proof}
If $G$ is a $p$-group and $Z_i \neq G$, then $G/Z_i$ is a nontrivial $p$-group, so
$Z_{i+1} \supsetneq Z_i$ by Proposition~\ref{groups:proposition-p-group-center}; hence
$Z_n = G$ eventually. The chain $Z_0 \trianglelefteq Z_1 \trianglelefteq \cdots$ has abelian
quotients, so nilpotent groups are solvable by
Proposition~\ref{groups:proposition-solvable}. $S_3$ is solvable via
$\{e\} \trianglelefteq A_3 \trianglelefteq S_3$, but $Z(S_3) = \{e\}$ (a non-identity permutation of
$\{1,2,3\}$ does not commute with every transposition), so $Z_i = \{e\}$ for all $i$.
\end{proof}

\begin{definition}\label{groups:definition-composition-series}
A \emph{composition series} of a group $G$ is a chain
$\{e\} = G_n \trianglelefteq G_{n-1} \trianglelefteq \cdots \trianglelefteq G_0 = G$ in which every factor
$G_i/G_{i+1}$ is simple; the factors are the \emph{composition factors} and $n$ is
the \emph{length}.
\end{definition}

\begin{theorem}[Jordan--H\"older]\label{groups:theorem-jordan-holder}
Every finite group has a composition series, and any two composition series of a
finite group have the same length and the same composition factors up to
isomorphism and permutation.
\end{theorem}

\begin{proof}
Existence: by induction on $|G|$; if $G \neq \{e\}$, a normal subgroup $N \neq G$ of
largest order has $G/N$ simple by the correspondence theorem, and $N$ has a
composition series.

Uniqueness, by induction on $|G|$. Let $(G_i)$ and $(H_j)$ be composition series of
lengths $n$ and $m$. If $G_1 = H_1$, apply induction to $G_1$. Otherwise
$H_1 \not\subseteq G_1$: if $H_1 \subsetneq G_1$, then $G_1/H_1$ would be a nontrivial proper
normal subgroup of the simple group $G/H_1$. So $G_1H_1$ is a normal subgroup strictly
containing $G_1$, and $G_1 H_1 = G$ because $G/G_1$ is simple.
Let $K = G_1 \cap H_1$. By Theorem~\ref{groups:theorem-isomorphism-theorems}(2),
$G_1/K \cong G/H_1$ and $H_1/K \cong G/G_1$, both simple. Choose a composition series of
$K$, of length $l$. Then $G_1 \trianglerighteq K \trianglerighteq \cdots$ is a composition series of $G_1$
of length $l + 1$; by induction applied to $G_1$ it has the same length and factors as
$G_1 \trianglerighteq G_2 \trianglerighteq \cdots$, so $n - 1 = l + 1$. Similarly $m - 1 = l + 1$. The factors of
$(G_i)$ are $G/G_1$, $G_1/K \cong G/H_1$ and those of $K$; the factors of $(H_j)$ are
$G/H_1$, $H_1/K \cong G/G_1$ and those of $K$. These agree up to permutation.
\end{proof}

\section{Free groups and presentations}\label{groups:section-free}

\begin{definition}\label{groups:definition-free-group}
Let $S$ be a set and $A = S \times \{1, -1\}$, whose elements we write $s$ and
$s^{-1}$; for $x = s^{\pm 1}$ write $x^{-1} = s^{\mp 1}$. A \emph{word} is a finite
sequence $x_1 \cdots x_k$ of elements of $A$, and it is \emph{reduced} if no two
adjacent letters are of the form $x x^{-1}$. Let $F(S)$ be the set of reduced words.
\end{definition}

\begin{theorem}\label{groups:theorem-free-group}
There is a unique group structure on $F(S)$ in which the product of reduced words
$u$ and $v$ is obtained from the concatenation $uv$ by repeatedly deleting adjacent
pairs $x x^{-1}$. With the map $\iota \colon S \to F(S)$ sending $s$ to the one-letter
word, $F(S)$ is the \emph{free group} on $S$: for every group $G$ and map
$f \colon S \to G$ there is a unique homomorphism $\bar f \colon F(S) \to G$ with
$\bar f \circ \iota = f$. So $F$ is left adjoint to the forgetful functor
$\cat{Grp} \to \Sets$.
\end{theorem}

\begin{proof}
For $x \in A$ define $L_x \colon F(S) \to F(S)$ by $L_x(w) = xw$ if $w$ does not begin
with $x^{-1}$, and $L_x(x^{-1}w') = w'$. Then $L_x L_{x^{-1}} = \id$ for all $x$, so each
$L_x$ is a bijection, and $L_{x^{-1}} = L_x^{-1}$. Let $\Phi \leq \mathrm{Sym}(F(S))$ be the
subgroup generated by the $L_x$, and $\epsilon \colon \Phi \to F(S)$, $\pi \mapsto \pi(\emptyset)$,
where $\emptyset$ is the empty word. By Lemma~\ref{groups:lemma-generated-subgroup}, every
$\pi \in \Phi$ is a product $L_{y_1} \cdots L_{y_m}$; deleting adjacent factors
$L_x L_{x^{-1}} = \id$ we may assume $y_1 \cdots y_m$ is reduced, and then
$\pi(\emptyset) = y_1 \cdots y_m$. Hence $\pi$ is determined by $\pi(\emptyset)$, and every reduced
word $w = x_1 \cdots x_k$ equals $L_{x_1} \cdots L_{x_k}(\emptyset)$. So $\epsilon$ is a bijection, and we
transport the group structure of $\Phi$ to $F(S)$. The product of $u$ and $v$ is then
$L_{u_1} \cdots L_{u_k}(v)$, which is computed by the deletion process of the
statement.

For $f \colon S \to G$ extend $f$ to $A$ by $f(s^{-1}) = f(s)^{-1}$ and define
$\bar f(x_1 \cdots x_k) = f(x_1) \cdots f(x_k)$. The value of $f(x_1)\cdots f(x_k)$ on an
arbitrary word does not change when an adjacent pair $x x^{-1}$ is deleted, and the
product in $F(S)$ is obtained from concatenation by such deletions; so $\bar f$ is a
homomorphism. It is unique because $\iota(S)$ generates $F(S)$.
\end{proof}

\begin{definition}\label{groups:definition-presentation}
Let $S$ be a set and $R \subseteq F(S)$. The group with \emph{presentation}
$\langle S \mid R \rangle$ is $F(S)/N$, where $N$ is the smallest normal subgroup containing $R$
(the intersection of all normal subgroups containing $R$). Every group $G$ has a
presentation: take $S = G$ and $R$ the kernel of $F(G) \to G$.
\end{definition}

\begin{proposition}\label{groups:proposition-presentation-property}
Let $G = \langle S \mid R \rangle$ and let $f \colon S \to H$ be a map to a group such that
$\bar f(r) = e$ for all $r \in R$. Then there is a unique homomorphism $G \to H$ sending
the class of each $s$ to $f(s)$.
\end{proposition}

\begin{proof}
$\ker \bar f$ is a normal subgroup containing $R$, hence containing $N$; apply
Theorem~\ref{groups:theorem-quotient}. Uniqueness holds because the classes of $S$
generate $G$.
\end{proof}

\begin{example}\label{groups:example-presentations}
\begin{enumerate}
\item $\langle a \mid a^n \rangle \cong \ZZ/n\ZZ$: the homomorphism $F(\{a\}) \to \ZZ/n\ZZ$,
$a \mapsto 1$, kills $a^n$ and so induces a surjection $\langle a \mid a^n \rangle \to \ZZ/n\ZZ$;
the source is generated by the class of $a$, which has order dividing $n$, so it has
at most $n$ elements.
\item $\langle a, b \mid aba^{-1}b^{-1} \rangle \cong \ZZ^2$: the map $a \mapsto (1, 0)$,
$b \mapsto (0, 1)$ induces a surjection, and a homomorphism
$\ZZ^2 \to \langle a, b \mid aba^{-1}b^{-1} \rangle$ inverse to it exists because the classes of $a$
and $b$ commute.
\item $D_n \cong \langle r, s \mid r^n, s^2, srsr \rangle$, by a similar argument.
\end{enumerate}
\end{example}

\begin{definition}\label{groups:definition-amalgam}
Let $\alpha_1 \colon H \to G_1$ and $\alpha_2 \colon H \to G_2$ be homomorphisms. The
\emph{amalgamated free product} $G_1 *_H G_2$ is the group with generators the
disjoint union of the sets $G_1$ and $G_2$ and relations $x y z^{-1}$ whenever $xy = z$
holds in $G_1$ or in $G_2$, and $\alpha_1(h)\alpha_2(h)^{-1}$ for $h \in H$. If $H$ is trivial,
it is the \emph{free product} $G_1 * G_2$.
\end{definition}

\begin{proposition}\label{groups:proposition-amalgam-pushout}
With the evident homomorphisms $G_i \to G_1 *_H G_2$, the amalgamated free product is a
pushout of $G_1 \leftarrow H \rightarrow G_2$ in $\cat{Grp}$. In particular $G_1 * G_2$ is a
coproduct, and $\cat{Grp}$ has all colimits.
\end{proposition}

\begin{proof}
The relations of the first kind make the maps $G_i \to G_1 *_H G_2$ homomorphisms, and
those of the second kind make the square commute. Given homomorphisms
$f_i \colon G_i \to K$ with $f_1\alpha_1 = f_2\alpha_2$, the map on generators given by $f_1$
and $f_2$ kills all relations, so it induces a homomorphism by
Proposition~\ref{groups:proposition-presentation-property}, which is unique since the
images of $G_1$ and $G_2$ generate. Coproducts of arbitrary families are
constructed in the same way, and coequalisers are quotients by normal subgroups;
so $\cat{Grp}$ has all colimits by the dual of
Theorem~\ref{categories:theorem-limits-products-equalisers}.
\end{proof}

\begin{remark}\label{groups:remark-normal-form}
If $\alpha_1$ and $\alpha_2$ are injective, the maps $G_i \to G_1 *_H G_2$ are injective,
and every element has a normal form as an alternating product of coset
representatives \cite[Chapter I, Theorem 1]{serre-trees},
\cite[Chapter IV]{lyndon-schupp}. For example
$\mathrm{SL}_2(\ZZ) \cong \ZZ/4\ZZ *_{\ZZ/2\ZZ} \ZZ/6\ZZ$ \cite[Chapter I, \S 4]{serre-trees}.
Amalgamated free products compute fundamental groups of unions of spaces via the
Seifert--van Kampen theorem (Chapter~\ref{homotopy}).
\end{remark}

\section{Abelian groups}\label{groups:section-abelian}

In this section groups are abelian and written additively. For $n \in \ZZ$ and
$a \in A$ we write $na$ for the $n$-th power of Lemma~\ref{groups:lemma-powers}.

\begin{definition}\label{groups:definition-free-abelian}
For a set $S$, let $\ZZ^{(S)}$ be the set of maps $S \to \ZZ$ that vanish outside a
finite set, with pointwise addition, and $e_s$ the map that is $1$ at $s$ and $0$
elsewhere. An abelian group is \emph{free} if it is isomorphic to some $\ZZ^{(S)}$,
and a family $(a_s)$ in $A$ is a \emph{basis} if the homomorphism
$\ZZ^{(S)} \to A$, $e_s \mapsto a_s$, is an isomorphism. The \emph{torsion subgroup} of $A$
is $A_{\mathrm{tors}} = \{a : na = 0 \text{ for some } n \geq 1\}$, and $A$ is
\emph{torsion-free} if $A_{\mathrm{tors}} = 0$.
\end{definition}

\begin{proposition}\label{groups:proposition-free-abelian}
For every abelian group $A$ and map $f \colon S \to A$ there is a unique homomorphism
$\ZZ^{(S)} \to A$ with $e_s \mapsto f(s)$, namely $\sum_s n_s e_s \mapsto \sum_s n_s f(s)$. If
$S$ is finite with $n$ elements, $\ZZ^{(S)} \cong \ZZ^n$, and $\ZZ^n \cong \ZZ^m$ implies $n = m$.
\end{proposition}

\begin{proof}
The formula defines a homomorphism, and it is forced. If $\ZZ^n \cong \ZZ^m$, then
$\ZZ^n/2\ZZ^n \cong \ZZ^m/2\ZZ^m$, and these groups have $2^n$ and $2^m$ elements.
\end{proof}

\begin{definition}\label{groups:definition-rank}
The \emph{rank} of a free abelian group $\ZZ^n$ is $n$; it is well defined by
Proposition~\ref{groups:proposition-free-abelian}.
\end{definition}

\begin{proposition}\label{groups:proposition-subgroups-free}
Every subgroup of $\ZZ^n$ is free of rank at most $n$.
\end{proposition}

\begin{proof}
By induction on $n$; for $n = 0$ it is clear. Let $H \leq \ZZ^n$ and
$p \colon \ZZ^n \to \ZZ$ the last coordinate. Then $H \cap \ker p$ is a subgroup of
$\ker p \cong \ZZ^{n-1}$, hence free of rank $\leq n - 1$ with some basis $(h_1, \ldots, h_k)$.
By Proposition~\ref{groups:proposition-subgroups-integers}, $p[H] = d\ZZ$. If $d = 0$,
$H = H \cap \ker p$. Otherwise choose $h \in H$ with $p(h) = d$; every $x \in H$ is uniquely
$x = (x - q h) + qh$ with $q = p(x)/d$ and $x - qh \in H \cap \ker p$, so $(h_1, \ldots, h_k, h)$
is a basis of $H$.
\end{proof}

\begin{theorem}[Smith normal form over $\ZZ$]\label{groups:theorem-smith}
Let $M$ be an $m \times n$ matrix with integer entries. There are invertible integer
matrices $P$ and $Q$ (with integer inverses) such that $PMQ$ is diagonal, with
diagonal entries $d_1, \ldots, d_r, 0, \ldots, 0$ where $d_i \geq 1$ and
$d_1 \mid d_2 \mid \cdots \mid d_r$.
\end{theorem}

\begin{proof}
Multiplying by invertible matrices on the left and right performs the elementary
operations: swapping two rows or columns, multiplying a row or column by $-1$, and
adding an integer multiple of one row (column) to another. We argue by induction on
$m + n$; if $M = 0$ there is nothing to do. Among all matrices obtainable from $M$ by
elementary operations, choose one, $M'$, whose smallest nonzero entry in absolute
value, $\delta$, is as small as possible; move it to position $(1,1)$ and make it
positive. Every entry $a$ of the first row or column is divisible by $\delta$: otherwise
subtracting a multiple of the first column (row) would produce the remainder of $a$
modulo $\delta$, a smaller nonzero entry. So we can clear the first row and column.
Every remaining entry $a$ is also divisible by $\delta$: otherwise adding its row to the
first row and then subtracting a multiple of the first column would create a smaller
nonzero entry. Apply induction to the remaining $(m-1) \times (n-1)$ block; the
operations on it do not affect divisibility of its entries by $\delta$, so $d_1 = \delta$
divides all later $d_i$.
\end{proof}

\begin{theorem}[Structure of finitely generated abelian groups]\label{groups:theorem-structure-abelian}
Let $A$ be a finitely generated abelian group. There are unique integers
$r \geq 0$ and $d_1, \ldots, d_k \geq 2$ with $d_1 \mid d_2 \mid \cdots \mid d_k$ such that
\[
A \cong \ZZ^r \oplus \ZZ/d_1\ZZ \oplus \cdots \oplus \ZZ/d_k\ZZ.
\]
Here $A_{\mathrm{tors}} \cong \bigoplus_i \ZZ/d_i\ZZ$, and $A/A_{\mathrm{tors}} \cong \ZZ^r$.
\end{theorem}

\begin{proof}
\emph{Existence.} Choose generators $a_1, \ldots, a_n$, giving a surjection
$\pi \colon \ZZ^n \to A$. By Proposition~\ref{groups:proposition-subgroups-free},
$\ker \pi$ is free with a basis $v_1, \ldots, v_m$, $m \leq n$. Let $M$ be the $n \times m$
matrix with columns $v_j$. Changing the basis of $\ZZ^n$ corresponds to multiplying $M$
by an invertible matrix on the left, and changing the basis of $\ker \pi$ to multiplying
on the right. By Theorem~\ref{groups:theorem-smith} there are bases
$f_1, \ldots, f_n$ of $\ZZ^n$ and $w_1, \ldots, w_m$ of $\ker \pi$ with $w_i = d_i f_i$ for
$i \leq s$ and $w_i = 0$ for $i > s$; since the $w_i$ are a basis, $s = m$ and all $d_i \geq 1$.
Hence $A \cong \ZZ^n / \ker\pi \cong \bigoplus_{i \leq m} \ZZ/d_i\ZZ \oplus \ZZ^{n-m}$, and we
discard the factors with $d_i = 1$.

\emph{Uniqueness.} In such a decomposition the torsion subgroup is
$T = \bigoplus \ZZ/d_i\ZZ$ and $A/T \cong \ZZ^r$, so $r$ is determined by
Proposition~\ref{groups:proposition-free-abelian}. For a prime $p$ and $j \geq 0$,
$p^j(\ZZ/d\ZZ)/p^{j+1}(\ZZ/d\ZZ)$ has order $p$ if $p^{j+1} \mid d$ and order $1$
otherwise. Hence $|p^jT/p^{j+1}T| = p^{t(p, j)}$, where $t(p, j)$ is the number of
$i$ with $p^{j+1} \mid d_i$. By the divisibility chain, these are the last $t(p,j)$ indices.
So the numbers $t(p, j)$, which depend only on $T$, determine the exponent of $p$ in
every $d_i$, provided $k$ is known; and $k = \max_p t(p, 0)$, since a prime dividing
$d_1$ divides all $d_i$.
\end{proof}

\begin{corollary}\label{groups:corollary-chinese-remainder}
If $\gcd(m, n) = 1$, then $\ZZ/mn\ZZ \cong \ZZ/m\ZZ \oplus \ZZ/n\ZZ$. Every finite abelian group is
a direct sum of cyclic groups of prime power order.
\end{corollary}

\begin{proof}
The element $(1, 1)$ of $\ZZ/m\ZZ \oplus \ZZ/n\ZZ$ has order $\operatorname{lcm}(m,n) = mn$, so it
generates. Apply this to the prime factorisations of the $d_i$.
\end{proof}

\begin{counterexample}\label{groups:counterexample-rationals-not-free}
The group $(\QQ, +)$ is torsion-free but not free, and not finitely generated. Any two
elements $a/b, c/d$ satisfy the nontrivial relation $bc \cdot (a/b) - ad \cdot (c/d) = 0$,
so a basis would have one element, and $\QQ \cong \ZZ$ is impossible since
$\QQ$ is divisible ($x = 2 \cdot (x/2)$) while $1 \in \ZZ$ is not twice anything. A
finitely generated torsion-free group would be free by
Theorem~\ref{groups:theorem-structure-abelian}.
\end{counterexample}

\section{Exercises}\label{groups:section-exercises}

\begin{exercise}\label{groups:exercise-burnside-lemma}
Let a finite group $G$ act on a finite set $X$. Show that the number of orbits is
$\frac{1}{|G|}\sum_{g \in G} |X^g|$, where $X^g = \{x : gx = x\}$.
\end{exercise}

\begin{solution}
Count the set $\{(g, x) : gx = x\}$ in two ways: it has $\sum_g |X^g|$ elements, and
also $\sum_x |G_x| = \sum_x |G|/|Gx|$. Each orbit $O$ contributes
$\sum_{x \in O} |G|/|O| = |G|$ to the last sum.
\end{solution}

\begin{exercise}\label{groups:exercise-small-orders}
Show that every group of order at most $5$ is abelian, and that $S_3$ is the smallest
non-abelian group.
\end{exercise}

\begin{exercise}\label{groups:exercise-abelianization-symmetric}
Show that $[S_n, S_n] = A_n$ for $n \geq 2$ and $S_n^{\mathrm{ab}} \cong \ZZ/2\ZZ$.
\end{exercise}

\begin{exercise}\label{groups:exercise-frattini}
(Frattini argument.) Let $N \trianglelefteq G$ be finite and $P$ a Sylow $p$-subgroup of $N$.
Show that $G = N \cdot N_G(P)$.
\end{exercise}

\begin{exercise}\label{groups:exercise-order-pq}
Let $p < q$ be primes with $p \nmid q - 1$. Show that every group of order $pq$ is cyclic.
\end{exercise}

\begin{exercise}\label{groups:exercise-free-group-rank}
Show that $F(S) \cong F(T)$ implies $|S| = |T|$ for finite sets $S$, $T$. (Hint: count
homomorphisms to $\ZZ/2\ZZ$.)
\end{exercise}

\begin{exercise}\label{groups:exercise-free-group-nonabelian}
Show that the free group on two generators is not solvable. (Hint: construct a
surjection onto $S_5$.)
\end{exercise}

\begin{exercise}\label{groups:exercise-q8}
Let $Q_8 = \{\pm 1, \pm i, \pm j, \pm k\}$ with $i^2 = j^2 = k^2 = ijk = -1$. Show that
every subgroup of $Q_8$ is normal but $Q_8$ is not abelian, and that $Q_8$ is not a
semidirect product of two proper subgroups.
\end{exercise}
